\section*{第 10 章}
\addcontentsline{toc}{section}{第 10 章}

\exercise{习题 10.1}
\textbf{题目：}
利用 $m$ 序列的移位相加特性，证明双极性 $m$ 序列的周期性自相关函数为二值函数，且主副峰之比等于码长（周期）。

\par\textbf{证明：}
设周期为 $P$ 的二进制 $m$ 序列为 $\{b_n\}$，其双极性形式记为
\[
a_n=(-1)^{b_n}\in\{+1,-1\}.
\]
周期性自相关函数定义为
\[
R_a(m)=\sum_{n=0}^{P-1}a_{n+m}a_n
=\sum_{n=0}^{P-1}(-1)^{b_{n+m}\oplus b_n}.
\]

当 $m\equiv 0\pmod P$ 时，
\[
b_{n+m}=b_n,
\]
故每项都为 $1$，于是
\[
R_a(0)=P.
\]

当 $m\not\equiv 0\pmod P$ 时，由 $m$ 序列的移位相加特性，序列
\[
\{b_{n+m}\oplus b_n\}
\]
仍是一个周期为 $P$ 的伪随机序列，其中 $1$ 的个数比 $0$ 的个数多 1。映射到双极性后，即 $-1$ 的个数比 $+1$ 的个数多 1，因此
\[
R_a(m)=-1,\qquad m\not\equiv 0\pmod P.
\]

所以双极性 $m$ 序列的周期性自相关函数只有两种取值：
\[
R_a(m)=
\begin{cases}
P, & m\equiv 0\pmod P,\\
-1, & m\not\equiv 0\pmod P.
\end{cases}
\]
故主峰与副峰之比为
\[
\frac{P}{1}=P,
\]
即等于码长（周期）。

\medskip
\exercise{习题 10.2}
\textbf{题目：}
已知线性反馈移存器序列的特征多项式为
\[
f(x)=x^3+x+1,
\]
求该序列的状态转移图，并说明它是否是 $m$ 序列。

\par\textbf{解：}
该序列发生器有 $3$ 个寄存器，递推关系可写为
\[
a_{n+3}=a_{n+1}+a_n.
\]
以 $(a_2,a_1,a_0)$ 为状态，不计全零状态，则状态循环为
\[
001\to 100\to 110\to 111\to 011\to 101\to 010\to 001.
\]
它经历了全部 $2^3-1=7$ 个非零状态，因此该序列的周期为 7，是\textbf{$m$ 序列}。

\medskip
\exercise{习题 10.3}
\textbf{题目：}
已知 $m$ 序列的特征多项式为
\[
f(x)=x^4+x+1,
\]
写出该序列一个周期中的所有游程。

\par\textbf{解：}
取初始状态为全 1，则可得一个周期序列为
\[
111101011001000.
\]
从中可读出共有 8 个游程：
\[
1111,\ 0,\ 1,\ 0,\ 11,\ 00,\ 1,\ 000.
\]
因此：
\[
\text{长度为 1 的游程有 }4\text{ 个},
\]
\[
\text{长度为 2 的游程有 }2\text{ 个},
\]
\[
\text{长度为 3 的游程有 }1\text{ 个},
\]
\[
\text{长度为 4 的游程有 }1\text{ 个}.
\]

\medskip
\exercise{习题 10.4}
\textbf{题目：}
已知优选对 $m_1,m_2$ 的特征多项式为
\[
f_1(x)=x^3+x+1,\qquad f_2(x)=x^3+x^2+1,
\]
写出由此优选对产生的所有 Gold 码，并求其中两个的周期互相关函数。

\par\textbf{解：}
由题给优选对可得两个 $m$ 序列分别为
\[
m_1=1110100,\qquad m_2=1110010.
\]
将 $m_2$ 作不同循环移位并与 $m_1$ 逐位模 2 相加，可得 7 个 Gold 码：
\[
0000110,\quad 0010001,\quad 0111111,\quad 1100011,
\]
\[
1011010,\quad 0101000,\quad 1001101.
\]

若取前两个 Gold 码作为一对，则其周期互相关函数取值为三值：
\[
R_{12}(m)=
\begin{cases}
-1, & k=0,1,3,5,\\
3, & k=2,\\
-5, & k=4.
\end{cases}
\]

\medskip
\exercise{习题 10.5}
\textbf{题目：}
采用 $m$ 序列测距，已知时钟频率为 $1\ \text{MHz}$，最远目标距离为 $3000\ \text{km}$，求 $m$ 序列长度（一周期的码片数）。

\par\textbf{解：}
最远目标的往返传播时延为
\[
T_{\max}=\frac{2\times 3000\times 10^3}{3\times 10^8}=0.02\ \text{s}=20\ \text{ms}.
\]
$m$ 序列一个周期的持续时间必须大于该最大时延。时钟频率为 $1\ \text{MHz}$，因此所需最少码片数为
\[
N>T_{\max}\times 10^6=2\times 10^4.
\]
满足该条件的最小 $m$ 序列长度为
\[
2^{15}-1=32767.
\]

\medskip
\exercise{习题 10.6}
\textbf{题目：}
已知某线性反馈移存器序列发生器的特征多项式为
\[
f(x)=x^3+x^2+1,
\]
画出结构图，写出它的输出序列（至少含一个周期），并指出其周期。

\par\textbf{解：}
由特征多项式
\[
f(x)=x^3+x^2+1
\]
可知反馈抽头在 3 级和 2 级，递推关系为
\[
a_{n+3}=a_{n+2}+a_n.
\]
故结构图为 3 级寄存器，第二级和第三级反馈到模 2 加法器后送回输入端。

取一组非零初始状态，可得其输出序列一个周期为
\[
1011100.
\]
因此该序列周期为
\[
P=7.
\]

\medskip
\exercise{习题 10.7}
\textbf{题目：}
\begin{enumerate}
\item 写出 8 阶哈达玛矩阵；
\item 验证其第 3 行和第 4 行正交。
\end{enumerate}

\par\textbf{解：}
8 阶哈达玛矩阵可写为
\[
H_8=
\begin{pmatrix}
1&1&1&1&1&1&1&1\\
1&-1&1&-1&1&-1&1&-1\\
1&1&-1&-1&1&1&-1&-1\\
1&-1&-1&1&1&-1&-1&1\\
1&1&1&1&-1&-1&-1&-1\\
1&-1&1&-1&-1&1&-1&1\\
1&1&-1&-1&-1&-1&1&1\\
1&-1&-1&1&-1&1&1&-1
\end{pmatrix}.
\]
其第 3 行与第 4 行内积为
\[
1\cdot 1+1\cdot(-1)+(-1)\cdot(-1)+(-1)\cdot 1+1\cdot 1+1\cdot(-1)+(-1)\cdot(-1)+(-1)\cdot 1=0.
\]
故第 3 行与第 4 行\textbf{正交}。

\medskip
\exercise{习题 10.8}
\textbf{题目：}
某直扩系统接收信号为
\[
s(t)=d(t)c(t)\cos(\omega_ct+\varphi)+c_0(t)\sin(\omega_ct+\varphi),
\]
其中 $d(t)\in\{\pm 1\}$ 为数据信号，$c(t),c_0(t)$ 为两路正交扩频码，$\varphi$ 为随机相位。设计相应接收框图。

\par\textbf{解：}
可采用\textbf{载波恢复 + 解扩判决}结构。一个可行方案为：
\begin{enumerate}
\item 对接收信号作 I/Q 解调，利用本地振荡器并结合 $c_0(t)$ 支路做相位跟踪，恢复随机相位 $\varphi$；
\item 载波同步后，取数据支路信号与本地扩频码 $c(t)$ 相乘解扩；
\item 在一个比特周期内积分
\[
\int_0^{T_b}(\cdot)\,dt
\]
并作判决，恢复 $d(t)$。
\end{enumerate}
也就是说，框图可描述为：
\[
s(t)\rightarrow \text{I/Q 解调}\rightarrow \text{相位跟踪}\rightarrow \times c(t)\rightarrow \int_0^{T_b}(\cdot)dt\rightarrow \text{判决}.
\]

\medskip
\exercise{习题 10.9}
\textbf{题目：}
某直扩系统接收信号为
\[
s(t)=d(t)c_1(t)\cos\omega_ct+d(t)c_2(t)\sin\omega_ct,
\]
其中 $d(t)\in\{\pm 1\}$ 为数据信号，$c_1(t),c_2(t)$ 为两路伪码扩频码，设计相应接收框图。

\par\textbf{解：}
可采用正交两支路分别解扩、再合并判决的结构：
\begin{enumerate}
\item 先对 $s(t)$ 作 I/Q 解调，得到同相支路和正交支路；
\item I 路与本地码 $c_1(t)$ 相乘，Q 路与本地码 $c_2(t)$ 相乘；
\item 两路输出相加后，在一个比特周期上积分；
\item 最后送入判决器恢复 $d(t)$。
\end{enumerate}
即框图可概括为
\[
s(t)\rightarrow \text{I/Q 解调}\rightarrow
\begin{cases}
\times c_1(t)\\
\times c_2(t)
\end{cases}
\rightarrow \Sigma \rightarrow \int_0^{T_b}(\cdot)dt\rightarrow \text{判决}.
\]

\medskip
